% ============================================================================
% Fichier  : partie5/P05_C21_anti_forensique.tex
% Titre    : Anti-Forensique et Contremesures
% Auteur   : MINKA MI NGUIDJOÏ Thierry Emmanuel
% Version  : 1.0 -- Juin 2026
% Statut   : RÉDIGÉ
% Sources  : Synthèse Parties II-IV (timestomping Ch.14, StopLogging Ch.17,
%            factory reset Ch.18, suppression $UsnJrnl Ch.14)
%            Affaires MVOM (suppression script), ONAMBELE
% Positionnement : Premier chapitre de la Partie V. Référencé en transition
%                  par Ch.14 (timestomping), Ch.17 (StopLogging CloudTrail),
%                  Ch.18 (factory reset IoT). Ce chapitre systématise : pour
%                  chaque technique d'effacement, la contre-détection
%                  correspondante. Structure miroir des Parties II-IV.
% ============================================================================

\chapter{Anti-Forensique et Contremesures}
\label{chap:anti-for}

\epigraph{%
    « L'absence de preuve est elle-même une preuve --- à condition
    de savoir où et comment cette absence a été créée. »%
}{--- Principe de l'anti-forensique inversée}

\niveauChapitre{Expert}{%
    Maîtrise complète de la Partie~IV (Ch.~\ref{chap:for-sys}
    à~\ref{chap:rapport}) indispensable. Ce chapitre suppose la
    connaissance des artefacts normaux pour en détecter l'absence
    anormale.%
}

\begin{boxobjectifs}
\begin{itemize}
    \item Comprendre la logique de l'anti-forensique~: pourquoi
          l'effacement parfait est techniquement très difficile.
    \item Identifier les techniques d'anti-forensique système~:
          timestomping, suppression sélective de logs, wiping.
    \item Détecter les tentatives d'anti-forensique réseau~:
          désactivation de logging, usurpation, chiffrement opportuniste.
    \item Reconnaître les techniques anti-forensique mobile et IoT~:
          factory reset, suppression d'historique, applications vault.
    \item Appliquer le principe de Locard inversé~: toute tentative
          d'effacement laisse elle-même une trace.
    \item Documenter une tentative d'anti-forensique dans le rapport
          forensique de façon juridiquement opposable.
\end{itemize}
\end{boxobjectifs}

\bigskip

% ============================================================================
\section*{La logique de l'anti-forensique}
% ============================================================================

Tout au long des Chapitres~\ref{chap:for-sys} à~\ref{chap:for-iot},
plusieurs scénarios ont montré des attaquants tentant d'effacer leurs
traces~: suppression du script \texttt{enc\_all.ps1} dans l'affaire MVOM
(Ch.~\ref{chap:for-sys}), \texttt{StopLogging} sur CloudTrail
(Ch.~\ref{chap:for-cloud}), tentatives de \textit{factory reset}
sur des appareils IoT (Ch.~\ref{chap:for-iot}). Dans chaque cas,
\textbf{la tentative a échoué à effacer toutes les traces}.

Ce n'est pas un hasard pédagogique~: c'est une réalité technique.
La redondance forensique (Ch.~\ref{chap:for-sys}, Tableau~14.1)
signifie qu'effacer \emph{toutes} les traces d'une action exige
de connaître et neutraliser \emph{chaque} artefact qu'elle a produit
--- une tâche d'une difficulté croissante avec la sophistication
des systèmes modernes.

\separateur

% ============================================================================
\section{Principes Généraux de l'Anti-Forensique}
\label{sec:principes-anti-for}
% ============================================================================

\subsection{Taxonomie des techniques anti-forensiques}

\begin{tableaucours}{p{3.0cm}p{3.5cm}p{5.0cm}}{%
    \tableauHeader{Catégorie} &
    \tableauHeader{Objectif} &
    \tableauHeader{Exemples}
}
\textbf{Destruction} & Éliminer la preuve définitivement &
    Wiping disque, suppression sécurisée, factory reset, destruction
    physique \\
\textbf{Obfuscation} & Rendre la preuve illisible ou non identifiable &
    Chiffrement, packing de malware, stéganographie, encodage \\
\textbf{Falsification} & Modifier la preuve pour induire en erreur &
    Timestomping, usurpation d'identité réseau, faux logs \\
\textbf{Minimisation} & Réduire la quantité de preuves produites &
    Désactivation de logging, mode furtif, utilisation de RAM uniquement \\
\textbf{Diversion} & Détourner l'attention de l'investigateur &
    False flag (faire accuser un tiers), bruit de fond généré
    artificiellement \\
\end{tableaucours}
\vspace{-0.8em}
\begin{center}{\small\itshape Tableau~21.1 --- Taxonomie des techniques anti-forensiques}\end{center}

\subsection{Le principe de Locard inversé}

\begin{boxdef}
\textbf{Principe de Locard inversé}\index{Locard inversé}~:
toute tentative d'effacement ou de modification d'une preuve produit
elle-même de nouvelles traces. L'acte de suppression est lui-même
un événement système qui laisse des artefacts~--- souvent dans des
emplacements que l'attaquant ne maîtrise pas ou ignore.
\end{boxdef}

Ce principe explique pourquoi, dans l'affaire MVOM, la suppression du
script \texttt{enc\_all.ps1} (Ch.~\ref{chap:for-sys}) a elle-même été
enregistrée dans le \texttt{\$UsnJrnl} avec horodatage précis~: l'acte
de \texttt{FileDelete} est un événement journalisé au même titre
que \texttt{FileCreate}.

\begin{boxCRO}
\textbf{Analyse CRO --- L'asymétrie fondamentale de l'anti-forensique}

$\mathcal{R}$ (côté attaquant)~: pour effacer parfaitement une trace,
l'attaquant doit connaître \emph{tous} les artefacts qu'elle a générés
--- une connaissance rarement complète, même pour un attaquant
sophistiqué.

$\mathcal{R}$ (côté investigateur)~: il suffit qu'\emph{un seul}
artefact survive parmi les nombreux générés (redondance forensique)
pour révéler l'activité initiale ET la tentative d'effacement.

$\mathcal{O}$~: une tentative d'effacement détectée renforce
paradoxalement l'opposabilité du dossier~: elle démontre l'intention
(\textit{mens rea}) de dissimuler une activité, un élément
souvent utile à la qualification pénale.

\cro{$\sim$}{$\downarrow$ pour l'attaquant~; $\uparrow\uparrow$ pour l'investigateur}{$\uparrow$ preuve d'intention}

\textbf{Conclusion~:} l'asymétrie structurelle favorise l'investigateur.
La sophistication nécessaire pour un effacement parfait dépasse les
moyens de la grande majorité des attaquants, y compris sophistiqués.
\end{boxCRO}

% ============================================================================
\section{Anti-Forensique Système~: Détection des Tentatives}
\label{sec:anti-for-systeme}
% ============================================================================

\subsection{Timestomping~: manipulation des horodatages}

Le Chapitre~\ref{chap:for-sys} a introduit le timestomping et sa
détection par croisement \texttt{\$MFT}/\texttt{\$UsnJrnl}. Ce chapitre
approfondit les techniques et leurs limites.

\begin{boxoutils}{Outils de timestomping et leurs limites}
\begin{lstlisting}[language=bash_forensique]
# Outils couramment utilises par les attaquants :
# - timestomp.exe (Metasploit)
# - SetMACE
# - PowerShell natif :
Set-ItemProperty -Path "C:\malware.exe" -Name CreationTime `
    -Value "2020-01-01 00:00:00"
Set-ItemProperty -Path "C:\malware.exe" -Name LastWriteTime `
    -Value "2020-01-01 00:00:00"

# CE QUE CES OUTILS NE PEUVENT PAS MODIFIER :
# 1. Le $UsnJrnl -- enregistre l'evenement de modification
#    des timestamps lui-meme (USN_REASON_BASIC_INFO_CHANGE)
MFTECmd.exe -f \$UsnJrnl:\$J --csv ./output/
# Chercher : Reason contient "BasicInfoChange"
# Cet evenement EST la preuve du timestomping

# 2. Le $LogFile (journal de transactions NTFS)
# Contient l'historique des transactions, y compris les
# modifications de timestamps, sur une fenetre plus courte
# (quelques heures a quelques jours selon l'activite disque)
LogFileParser.exe -f \$LogFile --csv ./output/logfile.csv

# 3. Les artefacts independants (Prefetch, UserAssist, ShimCache)
# ne sont PAS synchronises avec les modifications du $MFT
# -- ils gardent leurs propres horodatages independants
\end{lstlisting}
\end{boxoutils}

\begin{boxscenario}{Détection de timestomping --- méthode systématique}
\textbf{Étape~1~:} comparer les 4~horodatages MACE d'un fichier suspect.
Une cohérence parfaite (les~4~identiques à la seconde près) est
elle-même suspecte~--- un fichier créé normalement présente rarement
4~horodatages strictement identiques.

\textbf{Étape~2~:} croiser avec le \texttt{\$UsnJrnl}. Si le
\texttt{\$UsnJrnl} contient un événement \texttt{FileCreate} à une
date différente de celle affichée dans le \texttt{\$MFT}, c'est
une preuve directe de manipulation.

\textbf{Étape~3~:} croiser avec les artefacts indépendants
(Prefetch, UserAssist). Si Prefetch enregistre une exécution à
une date que le \texttt{\$MFT} affiche comme antérieure à la
création du fichier, c'est une incohérence logique irréfutable.

\textbf{Constatation type~:}
``Le \texttt{\$MFT} affiche pour le fichier \texttt{rootkit.dll}
une date de création du 12/08/2019. Le \texttt{\$UsnJrnl} enregistre
un événement \texttt{FileCreate} pour ce même fichier le 13/03/2026
à 02h19:15, soit 6~ans et 7~mois après la date affichée par le
\texttt{\$MFT}. Cette divergence constitue une preuve technique
de modification des horodatages (timestomping).''
\end{boxscenario}

\subsection{Suppression sélective de logs}

\begin{boxoutils}{Techniques de suppression de logs Windows et détection}
\begin{lstlisting}[language=bash_forensique]
# Technique attaquant : effacer le journal des evenements
wevtutil cl Security
wevtutil cl System
# Efface TOUT le journal -- action elle-meme detectable

# DETECTION : l'effacement génère un evenement
# Event ID 1102 : "The audit log was cleared"
# Cet evenement EST enregistre dans le journal Security
# AVANT que celui-ci soit vide, OU dans le journal qui suit

# Recherche de l'Event ID 1102 (preuve d'effacement de logs)
wevtutil qe Security /q:"*[System[(EventID=1102)]]" /f:text

# Technique plus sophistiquee : suppression selective d'evenements
# (modification directe du fichier .evtx)
# DETECTION : comparer le nombre d'evenements attendu
# (frequence normale) avec le nombre observe -- un "trou"
# dans la sequence des Event ID est suspect

# Verifier la continuite des RecordNumber (numero sequentiel)
# Si RecordNumber saute de 4521 a 4598, 76 evenements manquent
Get-WinEvent -Path security.evtx |
    Select-Object RecordId, TimeCreated, Id |
    Sort-Object RecordId
\end{lstlisting}
\end{boxoutils}

\begin{boxalert}
\textbf{L'Event~ID~1102~: la signature universelle de l'effacement Windows}

\texttt{wevtutil cl} (\textit{clear log}) est la commande la plus
utilisée par les attaquants pour effacer leurs traces. Elle génère
\emph{systématiquement} l'Event~ID~1102 (``Le journal d'audit a été
effacé''), qui survit à l'effacement lui-même car il est écrit
\emph{après} l'opération de nettoyage. \textbf{Chercher l'Event~ID~1102
doit être un réflexe systématique} dans toute investigation Windows
suspectant une dissimulation.
\end{boxalert}

\subsection{Wiping et suppression sécurisée}

\begin{boxoutils}{Détection de wiping de fichiers}
\begin{lstlisting}[language=bash_forensique]
# Outils de wiping courants : sdelete (Sysinternals), CCleaner,
# Eraser -- ecrivent des donnees aleatoires sur l'espace libere

# DETECTION 1 : signature de sdelete dans le $MFT
# sdelete renomme le fichier en une chaine de caracteres avant
# suppression -- cette sequence de renommages est visible dans
# le $UsnJrnl
MFTECmd.exe -f \$UsnJrnl:\$J --csv ./output/ |
    grep -E "RenameOldName|RenameNewName"
# Plusieurs renommages successifs rapides = signature sdelete

# DETECTION 2 : entropie de l'espace libere
# Wiping remplit l'espace de donnees aleatoires (haute entropie)
# Espace libere normalement (juste marque disponible) garde
# le contenu original jusqu'a reecriture -- entropie variable
photorec /d /output/ -cmd /dev/sdb1 freespace
# Analyser l'entropie de l'espace libere :
ent /output/freespace_dump.bin
# Entropie proche de 8.0 bits/byte sur de larges plages = wiping

# DETECTION 3 : logs d'execution de l'outil de wiping lui-meme
# (ironie : l'outil de suppression de preuves laisse ses propres
# traces -- Prefetch, UserAssist, ShimCache pour sdelete.exe)
PECmd.exe -f SDELETE.EXE-XXXXXXXX.pf
\end{lstlisting}
\end{boxoutils}

% ============================================================================
\section{Anti-Forensique Réseau~: Désactivation et Camouflage}
\label{sec:anti-for-reseau}
% ============================================================================

\subsection{Désactivation de logging~: le cas CloudTrail}

Le Chapitre~\ref{chap:for-cloud} a présenté la tentative
\texttt{StopLogging} dans l'affaire MVOM. Ce mécanisme se généralise~:

\begin{tableaucours}{p{3.0cm}p{4.0cm}p{4.5cm}}{%
    \tableauHeader{Plateforme} &
    \tableauHeader{Action de désactivation} &
    \tableauHeader{Contre-mesure architecturale}
}
AWS CloudTrail & \texttt{StopLogging}, \texttt{DeleteTrail} &
    Object Lock (WORM) sur le bucket~S3~; alerte CloudWatch
    sur ces actions \\
Azure Activity Log & Désactivation via portail (rare, nécessite
    droits élevés) & Logs exportés vers Log~Analytics
    (rétention indépendante) \\
Logs Windows & \texttt{wevtutil cl} &
    Forwarding vers serveur Syslog centralisé (SIEM) en temps réel \\
journald (Linux) & \texttt{journalctl --vacuum-time} &
    Forwarding rsyslog vers serveur distant~; \texttt{Storage=persistent} \\
\end{tableaucours}
\vspace{-0.8em}
\begin{center}{\small\itshape Tableau~21.2 --- Désactivation de logging et contre-mesures}\end{center}

\begin{boiteRemarque}{Le principe du forwarding~: la meilleure défense}
La contre-mesure la plus efficace contre toute désactivation de
logging \emph{locale} est le \textbf{forwarding en temps réel} vers
un système distant que l'attaquant ne contrôle pas. Un attaquant
qui compromet un serveur et désactive ses logs locaux ne peut pas
effacer ce qui a déjà été transmis ailleurs. Cette architecture
(SIEM centralisé) devrait être recommandée dans toute mission de
sécurisation post-incident.
\end{boiteRemarque}

\subsection{Chiffrement opportuniste et tunneling}

\begin{boxoutils}{Détection de canaux chiffrés non identifiés}
\begin{lstlisting}[language=bash_forensique]
# L'attaquant chiffre son C2 pour eviter l'inspection du contenu
# (vu en Ch. 15 : tunneling DNS, HTTPS legitime detourne)

# Detection par metadonnees (meme si contenu illisible) :
# 1. Certificat TLS auto-signe ou recent
tshark -r capture.pcap -Y "tls.handshake.type == 11" \
    -T fields -e tls.handshake.certificate
# Verifier la date d'emission du certificat -- recent = suspect

# 2. JA3/JA3S fingerprinting (signature du client TLS)
# Outil : ja3 (python)
python3 ja3.py capture.pcap
# Comparer avec des bases de signatures malveillantes connues
# (ja3er.com, abuse.ch SSL Blacklist)

# 3. Pattern de connexion meme sans dechiffrement
# (beaconing -- vu en Ch. 15 -- fonctionne meme sur trafic chiffre,
# car les METADONNEES de connexion restent visibles)
\end{lstlisting}
\end{boxoutils}

% ============================================================================
\section{Anti-Forensique Mobile et IoT}
\label{sec:anti-for-mobile-iot}
% ============================================================================

\subsection{Factory reset~: ce qui survit réellement}

Le Chapitre~\ref{chap:for-iot} a introduit la limite du factory reset.
Cette section approfondit, avec application au mobile (Ch.~\ref{chap:for-mob}).

\begin{tableaucours}{p{3.2cm}p{4.0cm}p{4.5cm}}{%
    \tableauHeader{Données} &
    \tableauHeader{Effacées par factory reset~?} &
    \tableauHeader{Source de récupération}
}
Stockage interne (\texttt{/data}) & Oui (en général) &
    Carving sur l'espace non réécrit (extraction physique avant
    réécriture complète) \\
Carte SD externe & \textbf{Non} (sauf reformatage explicite) &
    Extraction directe de la SD \\
Logs opérateur (appels, SMS, data) & Non (jamais sur l'appareil) &
    Réquisition opérateur (art.~57 Loi~2010) \\
Données cloud synchronisées
    (avant le reset) & Non &
    Réquisition fournisseur (Google, iCloud) \\
Historique réseau opérateur
    (cellules GSM) & Non &
    Réquisition + analyse Cell~ID (Ch.~\ref{chap:for-mob}) \\
\end{tableaucours}
\vspace{-0.8em}
\begin{center}{\small\itshape Tableau~21.3 --- Factory reset~: ce qui est réellement effacé}\end{center}

\subsection{Applications ``vault'' et conteneurs cachés}

\begin{boxalert}
\textbf{Applications vault~: dissimulation par camouflage}

Les applications de type ``vault'' (\textit{Calculator+},
\textit{Vault -- Hide Photos}) se présentent sous l'apparence
d'une application anodine (calculatrice) mais cachent un conteneur
chiffré accessible par un code secret.

\textbf{Détection~:}
\begin{itemize}
    \item \texttt{adb shell pm list packages -f} révèle le nom de
          package réel, souvent différent du nom affiché
          (ex.~: \texttt{com.vault.hideit} pour une icône
          ``Calculatrice'')~;
    \item Analyse de la taille du stockage de l'application
          (\texttt{adb shell dumpsys package <package>}) -- une
          calculatrice de 200~Mo est suspecte~;
    \item Liste des permissions demandées (accès photos, stockage)
          incohérente avec la fonction affichée.
\end{itemize}
\end{boxalert}

\subsection{Suppression sélective WhatsApp/Telegram}

\begin{boxoutils}{Détection de suppression de conversations}
\begin{lstlisting}[language=bash_forensique]
-- Sur msgstore.db (WhatsApp, Ch. 16), verifier les "trous"
-- dans la sequence des _id (cles primaires auto-incrementees)
SELECT _id, timestamp FROM messages ORDER BY _id;
-- Si _id saute de 1847 a 1901, 53 messages ont ete supprimes

-- WhatsApp conserve des backups automatiques quotidiens
-- (msgstore-YYYY-MM-DD.1.db.crypt15) sur 7 jours par defaut
-- Comparer le backup de la veille avec la base actuelle :
-- les messages presents dans le backup mais absents de la base
-- actuelle ont ete supprimes par l'utilisateur

-- Telegram (cache local, moins de retention) :
-- chercher les fichiers cache/ supprimes via carving
-- (Ch. 9 -- foremost/scalpel) sur l'extraction du telephone
\end{lstlisting}
\end{boxoutils}

% ============================================================================
\section{Documenter une Tentative d'Anti-Forensique dans le Rapport}
\label{sec:documentation-anti-for}
% ============================================================================

La détection d'une tentative d'anti-forensique doit être documentée
selon les mêmes règles que toute autre constatation
(Ch.~\ref{chap:rapport})~--- avec une attention particulière car
elle constitue souvent un élément d'intention (\textit{mens rea}).

\begin{lstlisting}[language=bash_forensique,
    caption={Exemple de constatation et analyse --- tentative d'anti-forensique}]
CONSTATATIONS (anti-forensique) :

C.10 : Le journal des événements Security du serveur enregistre
       un événement Event ID 1102 ("Le journal d'audit a été
       effacé") à 02h19:48, par le compte "admin".
       [Source : wevtutil, Annexe H]

C.11 : Le $UsnJrnl enregistre la suppression du fichier
       "enc_all.ps1" à 02h18:30, 1 minute et 18 secondes avant
       l'effacement du journal Security (C.10).
       [Source : MFTECmd, Annexe B]

ANALYSE :

La séquence temporelle C.10-C.11 est cohérente, avec un niveau
de confiance élevé, avec une tentative de dissimulation des
traces de l'intrusion par l'auteur de l'attaque. Cette tentative
n'a cependant pas atteint son objectif : le $UsnJrnl, distinct
du journal des événements Security, conserve l'historique complet
des opérations sur le système de fichiers, incluant la création
et la suppression du script malveillant.

Hypothèse alternative envisagée et écartée : une maintenance
système légitime aurait pu causer cet effacement de logs.
Cette hypothèse est écartée car (1) aucune fenêtre de maintenance
n'était planifiée à cette date selon les procédures de
l'entreprise (Annexe E) et (2) la proximité temporelle immédiate
avec la suppression du fichier malveillant (78 secondes) rend
cette coïncidence hautement improbable.
\end{lstlisting}

\begin{boxjuris}
\textbf{Valeur juridique de la tentative d'anti-forensique}

En droit camerounais, la tentative de dissimulation de preuves
peut constituer un élément aggravant ou un indice de l'intention
coupable. L'art.~88 \loicam{2010/012} (refus de remise de clé de
déchiffrement, Ch.~\ref{chap:droit-cam}) illustre ce principe~:
le législateur sanctionne spécifiquement les comportements
d'obstruction à l'investigation. La documentation rigoureuse
d'une tentative d'effacement (avec ses propres constatations
et sources) renforce la valeur probatoire du dossier dans
son ensemble.
\end{boxjuris}

\separateur

\begin{boxresume}
\textbf{Ce qu'il faut retenir de ce chapitre~:}
\begin{itemize}
    \item \textbf{Principe de Locard inversé}~: toute tentative
          d'effacement produit elle-même de nouvelles traces.
          L'asymétrie favorise structurellement l'investigateur~:
          un seul artefact survivant suffit à révéler l'activité
          \emph{et} la tentative de dissimulation.

    \item \textbf{Timestomping}~: détecté par croisement
          \texttt{\$MFT}/\texttt{\$UsnJrnl}~; les outils de
          timestomping ne peuvent pas modifier le \texttt{\$UsnJrnl}
          ni les artefacts indépendants (Prefetch, UserAssist).

    \item \textbf{Event~ID~1102}~: signature universelle de
          l'effacement des journaux Windows (\texttt{wevtutil cl}).
          Réflexe systématique de recherche dans toute investigation
          suspectant une dissimulation.

    \item \textbf{Le forwarding en temps réel} (SIEM centralisé)
          est la meilleure architecture de défense~: un attaquant
          qui efface les logs locaux ne peut rien contre ce qui a
          déjà été transmis ailleurs.

    \item \textbf{Factory reset}~: n'efface ni la SD~Card externe,
          ni les logs opérateur, ni les données cloud déjà
          synchronisées. Toujours compléter par réquisition.

    \item \textbf{Applications vault}~: détectées par le nom de
          package réel (\texttt{pm list packages -f}), la taille
          de stockage incohérente, les permissions excessives.

    \item \textbf{Documentation}~: une tentative d'anti-forensique
          se documente comme toute constatation (source, horodatage)
          suivie d'une analyse avec niveau de confiance et
          hypothèses alternatives réfutées (E9bis PIU). Elle
          constitue souvent un indice d'intention coupable.
\end{itemize}
\end{boxresume}

\bigskip

\begin{boxexo}[title={\ding{45}\quad Exercice~21.1 --- Identifier les techniques \niveaubase}]
Pour chacune des observations suivantes, identifiez la catégorie
d'anti-forensique (Tableau~21.1) et la contre-mesure de détection
appropriée~:
\begin{enumerate}[(a)]
    \item Un fichier affiche une date de création de 2015 mais
          son contenu fait référence à un événement de 2026.
    \item Le journal Security d'un serveur Windows est vide pour
          toute la journée du 13~mars~2026.
    \item Un suspect possède une application ``Calculatrice''
          qui pèse 340~Mo sur son téléphone.
    \item Un trafic réseau chiffré présente des connexions toutes
          les exactement 300~secondes vers la même IP.
\end{enumerate}
\end{boxexo}

\begin{boxexo}[title={\ding{45}\quad Exercice~21.2 --- Détection de timestomping \niveaumoyen}]
L'analyse \texttt{\$MFT} d'un fichier \texttt{backdoor.dll} montre~:
\begin{lstlisting}[language=bash_forensique]
Created (E):   2018-03-22 09:15:00
Modified (M):  2018-03-22 09:15:00
Accessed (A):  2018-03-22 09:15:00
Changed (C):   2018-03-22 09:15:00
\end{lstlisting}
L'analyse \texttt{\$UsnJrnl} montre~:
\begin{lstlisting}[language=bash_forensique]
2026-03-13 02:19:15  backdoor.dll  FileCreate
2026-03-13 02:19:16  backdoor.dll  DataExtend|Close
2026-03-13 02:19:20  backdoor.dll  BasicInfoChange
\end{lstlisting}
\begin{enumerate}[(a)]
    \item Interprétez chacune des trois lignes du \texttt{\$UsnJrnl}.
          Que signifie spécifiquement \texttt{BasicInfoChange}~?
    \item Rédigez la constatation forensique correspondante,
          en respectant la distinction constatation/analyse.
    \item Pourquoi les 4~horodatages MACE parfaitement identiques
          dans le \texttt{\$MFT} sont-ils eux-mêmes un indice
          suspect, indépendamment de la divergence avec le
          \texttt{\$UsnJrnl}~?
\end{enumerate}
\end{boxexo}

\begin{boxexo}[title={\ding{45}\quad Exercice~21.3 --- Investigation complète avec anti-forensique \niveauexpert}]
Dans l'affaire ONAMBELE (Ch.~\ref{chap:gr-affaires}), vous découvrez
les éléments suivants lors de l'analyse du serveur compromis~:
\begin{itemize}
    \item Event~ID~1102 dans le journal Security à 03h42:15~;
    \item \texttt{\$UsnJrnl}~: suppression de 14~fichiers \texttt{.xlsx}
          entre 03h40:00 et 03h41:30~;
    \item Un fichier \texttt{cleanup.bat} dans le Prefetch, exécuté
          à 03h39:50, contenant (reconstitué depuis la RAM via
          Volatility \texttt{cmdline})~:
\begin{lstlisting}[language=bash_forensique]
del C:\Finance\rapports\*.xlsx
wevtutil cl Security
wevtutil cl System
sdelete -p 3 C:\Finance\rapports\*
\end{lstlisting}
\end{itemize}
\begin{enumerate}[(a)]
    \item Reconstituez la chronologie complète des événements,
          en distinguant ce qui s'est passé et ce qui a été
          tenté pour le dissimuler.
    \item L'utilisation de \texttt{sdelete -p~3} (3~passes
          d'écrasement) rend-elle les fichiers \texttt{.xlsx}
          définitivement irrécupérables~? Quelles sources
          alternatives existent malgré cela (pensez aux chapitres
          précédents~: sauvegardes, logs réseau, copies cloud)~?
    \item Rédigez 4~constatations et l'analyse correspondante,
          en respectant scrupuleusement la structure du
          Ch.~\ref{chap:rapport}.
    \item Quelle valeur cette séquence (suppression~$\to$
          effacement logs~$\to$ wiping sécurisé) apporte-t-elle
          à la démonstration de l'intention coupable~? Comment
          formuler cela sans franchir la limite de la qualification
          juridique (réservée à l'autorité judiciaire)~?
\end{enumerate}
\end{boxexo}

\bigskip

\begin{boxinfo}
\textbf{Transition.} La maîtrise de l'anti-forensique complète le
socle technique du cours~: vous savez désormais non seulement
investiguer, mais aussi reconnaître et contrer les tentatives de
dissimulation. Cette compétence est transversale~--- elle
s'applique à chaque type d'evidence rencontré dans la Partie~IV.

Les chapitres suivants de la Partie~V (selon le programme retenu)
abordent les perspectives prospectives~: OSINT avancé, intelligence
artificielle appliquée à la détection d'anomalies, et impact
des technologies post-quantiques sur l'intégrité des preuves
numériques à long terme.
\end{boxinfo}